\documentclass[11pt]{scrartcl}
\usepackage[usenames,dvipsnames,svgnames]{xcolor}
\usepackage[framemethod=TikZ]{mdframed}
\usepackage{amsmath,amssymb,amsthm,thmtools}

\mdfdefinestyle{bluebox}{roundcorner=10pt, linecolor=blue, backgroundcolor=TealBlue!5}
\declaretheoremstyle[headfont=\sffamily\bfseries\color{MidnightBlue}, mdframed={style=bluebox}, headpunct={\\[3pt]}]{thmblue}
\mdfdefinestyle{greenbar}{linecolor=ForestGreen, backgroundcolor=ForestGreen!5, linewidth=2pt, rightline=false, topline=false, bottomline=false}
\declaretheoremstyle[headfont=\bfseries\sffamily\color{ForestGreen!70!black}, mdframed={style=greenbar}, headpunct={ --- }]{thmgreen}
\declaretheorem[style=thmblue,name=Theorem,numberwithin=subsection]{theorem}
\declaretheorem[style=thmblue,name=Lemma,sibling=theorem]{lemma}
\declaretheorem[style=thmblue,name=Theorem,numbered=no]{theorem*}
\declaretheorem[style=thmgreen,name=Definition,sibling=theorem]{definition}
\declaretheorem[style=thmblue,name=Problem,numberwithin=subsection]{problem}

\begin{document}
\section{Sets}
\subsection{Basics}
\begin{definition}\label{def:set}
A \emph{set} is a collection of objects.
\end{definition}
\begin{theorem}[Cantor]\label{thm:cantor}
No set surjects onto its power set.
\end{theorem}
\begin{lemma} A helper. \end{lemma}
\begin{theorem*} An unnumbered statement. \end{theorem*}
\begin{problem}[5 points]\label{prob:one}
Use Theorem~\ref{thm:cantor} and Definition~\ref{def:set}.
\end{problem}
\iffalse
\begin{theorem} Hidden. \end{theorem}
\fi
\subsection{More}
\begin{problem} Second subsection. \end{problem}
\end{document}
